\subsection{Boundary CFT} 

I will first give a very brief review of fundamental concepts and tools in Boundary Conformal Field Theory (BCFT). For a more detailed introduction, one can refer to my CFT and BCFT note, click \href{https://xiao-daihua.github.io/note.html}{here}, and the following references \cite{cardyBoundaryConditionsFusion1989, northeYoungResearchersSchool2025, recknagelBoundaryConformalField}.

\subsubsection{Gluing Condition}

A Boundary CFT (BCFT) is a conformal field theory on a manifold with boundary induced from a parent CFT. It inherets the operator algebra data of the parent CFT, while having extra data due to the presence of boundaries. 

The simplest case is to consider a BCFT on the upper half plane (UHP) $\{z=x+iy|y\geq 0\}$, where the real axis $y=0$ serves as the boundary. 
To preserve conformal symmetry at the boundary, we need to impose the gluing condition on the energy-momentum tensor $T(z)$ and $\bar{T}(\bar{z})$ at the boundary:
\defi{
  \textbf{Gluing Condition}

  At the boundary $y=0$, the energy-momentum tensor satisfies:
  \begin{align}
    T(z) = \bar{T}(\bar{z}) \quad \text{at } y=0.
  \end{align}
}
With this gluing condition, we can prove that the conformal symmetry is "partially" preserved, which leads to the existence of a single copy of Virasoro algebra as symmetry algebra instead of two:
\thm{
  \textbf{Symmetry Algebra of BCFT}

  In a BCFT defined on the upper half plane with the gluing condition $T(z) = \bar{T}(\bar{z})$ at the boundary $y=0$, the symmetry algebra is generated by a single Virasoro algebra with generators:
  \begin{align}
    L_n^H = \frac{1}{2\pi i} \oint_{C} dz \, z^{n+1} T^H(z) 
  \end{align}
  where $ T^H(z) $ is the analytical continuation of $ T(z) $ and $ \bar{T}(\bar{z}) $. The generators satisfy:
  \begin{align}
    [L_n^H, L_m^H] = (n-m) L_{n+m}^H + \frac{c}{12} n(n^2-1) \delta_{n+m,0}
  \end{align}
}

\subsubsection{Boundary States Formalism}

There exist various ways of defining a BCFT from a parent CFT on a manifold by imposing different types of boundary conditions at the boundary. A tool of labeling and classifying these boundary conditions is the boundary state formalism. We define the boundary states as;
\defi{
  \textbf{Boundary State}

  Boundary states $\ket{\alpha}$ and $\ket{\beta}$ are defined as states living on
  the boundary of the parent CFT placed on a chopped cylinder, and they satisfy:
  \begin{itemize}
    \item The cylinder has the same geometry as a compactified BCFT on a strip,
          with length $L$ and circumference $\beta$.
    \item The partition function of the parent CFT on a cylinder with boundary
          states $\ket{\alpha}$ and $\ket{\beta}$ at its two ends is equivalent to
          the partition function of the BCFT on a strip with boundary conditions
          $\alpha$ and $\beta$ at the two ends. Mathematically:
  \end{itemize}
  \begin{align}
    Z_{\alpha\beta} = \bra{\alpha}\,
           \tilde{q}^{\frac{1}{2}(L_0 + \bar{L}_0 - c/12)}\,
         \ket{\beta} = \operatorname{Tr}_{\mathcal{H}_{\alpha\beta}}
      \left( q^{\,L_0^{H} - c/24} \right),
  \end{align}
  where $q = e^{-\pi \beta / L}$ and $\tilde{q} = e^{-4\pi L / \beta}$.Here $L_0$ and $\bar{L}_0$ are the Virasoro generators of the parent CFT, while $L_0^{H}$ is the Virasoro generator in the BCFT. These condition are often referred to the Cardy's condition.
}
We can use the boundary states as a state in the Hilbert space of the parent CFT to encode the boundary conditions of the BCFT. To be consistent with the gluing condition, the boundary states must satisfy the following Ishibashi condition:
\thm{
  \textbf{Ishibashi Condition}

  Boundary states $\ket{\alpha}$ should satisfy:
  \begin{align}
    (L_n - \bar{L}_{-n}) \ket{\alpha} = 0, \quad \forall n \in \mathbb{Z}.
  \end{align}
}
A solution to the Ishibashi condition can be constructed from each primary state $\ket{h}$ of the parent CFT, leading to the definition of Ishibashi states:
\defi{
  \textbf{Ishibashi State}

  For each primary state $\ket{h}$ of the parent CFT, the corresponding Ishibashi state $\ket{h}\rangle$ is defined as:
  \begin{align}
    \ket{h}\rangle = \sum_{N} \ket{h;N} \otimes \overline{\ket{h;N}},
  \end{align}
  where $\{\ket{h;N}\}$ is an orthonormal basis of the Verma module built on the chiral primary state $\ket{h}$. 
}

\subsubsection{Cardy's Condition}

Consider the parent CFT to be a diagonal CFT. Indeed, the main subject of this note, the Ising CFT, is a diagonal CFT. There exists an equivalent formulation of Cardy’s condition:
\thm{
  \textbf{Cardy's Condition}(Diagonal CFT)

  For two boundary states $ \ket{\alpha} $ and $ \beta $ they should satisfy that:
  \begin{align}
    \sum_h\langle\alpha|h\rangle\rangle\langle\langle h|\beta\rangle S_h^{h^{\prime}}=n_{\alpha\beta}^{h^{\prime}} \quad \text{or equivalently} \quad  \langle a|h^{\prime}\rangle\rangle\langle\langle h^{\prime}|b\rangle=\sum_hS_h^{h^{\prime}}n_{ab}^h.
  \end{align}
  where $ S^{h'}_h $ is the modular S matrix of the parent CFT and $ n_{\alpha\beta}^{h^{\prime}} $ are non-negative integers representing the multiplicities of the representation with highest weight $ h' $ in the open channel partition function $ Z_{\alpha\beta} $.
}
% If we assume that the boundary states are linear combination of Ishibashi States, we can also write the Cardy's condition a relation of expansion coefficients:
% \thm{\label{thm:cardycondition}
%   \textbf{Cardy's Condition (coefficient constrain)}
%
%   Consider the boundary states written as:
%   \begin{align}
%     \left|a\right\rangle=\sum_{i}B_{i}^{a}\left|i\right\rangle\rangle
%   \end{align}
%   The Cardy's condition says that:
%   \begin{align}
%     \sum_{ij} B_i^{a *}B_i^bS_{i}^j\chi_j(q)=\sum_in_{ab}^i\chi_i(q)\quad \Rightarrow \quad \sum_{i} B_i^{a *}B_i^bS_{i}^j=n_{ab}^j
%   \end{align}
%   where $ \chi_i $ is the character of the Virasoro representation. 
% }

A solution to the Cardy's condition is given by comparing it with the Verlinde formula. These solutions are called the Cardy States:
\defi{
  \textbf{Cardy States}

  Cardy states are linear combination of Ishibashi states with coefficients given by:
  \begin{align}
    |a\rangle=\sum_{j}\frac{S_{a}^{j}}{(S_{0}^{j})^{1/2}}|j\rangle\rangle
  \end{align}
}
Proof: Taking $ n_{ab}^i $ as the fusion matrix $ N^a_{bc} $ and Cardy states into Cardy's condition, we can see that the Cardy's condition exactly gives out the Verlinde formula.

\subsubsection{Elementary Boundary States}

Looking at the Cardy's condition we might see that, if a state $ \ket{a} $ satisfy the condition then $ \mathbb{N}_+\ket{a} $ must also satisfy this condition, for we only insist $ n^i_{ab} $ to be non-negative integers. 
However, we usually only consider the elementary boundary states that can not be decomposed into a sum of other boundary states. Thus we have the definition:
\defi{
  \textbf{Elementary Boundary States}

  For a parent CFT, a induced BCFT have following elementary boundary states $ \{\ket{a_i}\} $ that shall satisfy the following normalization and orthogonality condition: 
  \begin{align}
    n_{a_ia_j}^0 = \delta_{ij}
  \end{align}
}
The normalization is physical, this is because for a field theory, we assume to have a unique vaccum state, thus if $ n_{ab}^0 \geq 2 $ we will have multiple vacuum states in the open channel, which is unphysical, and shall be interpret as putting two theory togethor. Our goal is to find all elementary boundary states of a BCFT induced from a parent CFT.
Moreover, we can see that:
\lmm{
 Cardy states are elementary boundary states. 
}
This is because, cardy's state is a solution when considering $ n_{ab}^i = N_{ab}^i $ which is the bulk fusion matrix. And we know that according to the 2 point function of bulk primary fields, the vacuum representation only appear once in the fusion of two primary fields when they are the same, written in operator algebra as:
\begin{align}
  \phi_a \times \phi_b = \mathbb{1} \delta_{ab} + \cdots
\end{align}

\subsection{Ising CFT}

Ising CFT as the simplist unitary minimal model $ M(3/4) $ with central charge $ c=1/2 $. Here I will briefly review some basic data of Ising CFT that will be used in the following sections. For a more detailed introduction to Ising CFT, one can refer to \cite{difrancescoConformalFieldTheory1997}.

\subsubsection{Ising CFT Data}
Ising CFT is the simplest A-series unitary minimal model $M(3,4)$ with central charge $c=\tfrac{1}{2}$. 
It has three primary fields:
\begin{itemize}
  \item Identity operator $\mathbb{1}$ with conformal weights 
    $h_{\mathbb{1}}=\bar{h}_{\mathbb{1}}=0$, 
  \item Spin operator $\sigma$ with conformal weights 
    $h_{\sigma}=\bar{h}_{\sigma}= \tfrac{1}{16}$
  \item Energy operator $\epsilon$ with conformal weights 
    $h_{\epsilon}=\bar{h}_{\epsilon}=\tfrac{1}{2}$.
\end{itemize}
They satisfy the folowing operator algebra (fusion rules):
\begin{align}
  &\sigma \times \sigma = \mathbb{1} + \epsilon, \\
  &\sigma \times \epsilon = \sigma, \\
  &\epsilon \times \epsilon = \mathbb{1}.
\end{align}

\subsubsection{Virasoro Characters of Ising CFT}

Charaters of a Virasoro Representation is defined as :
\defi{
  \textbf{Character}

The character $\chi_h(q)$ of a Virasoro representation with highest weight $h$ is defined as:
\begin{align}
  \chi_h(q) = \operatorname{Tr}_{\mathcal{H}_h} \left( q^{L_0 - c/24} \right),
\end{align}
}
For Ising CFT the characters of the three primary fields are given by:
\begin{align}
  &\chi_0(q)=\frac{1}{2}q^{-\frac{1}{48}}\left(\prod_{n=0}^\infty\left(1+q^{\frac{1}{2}+n}\right)+\prod_{n=0}^\infty\left(1-q^{\frac{1}{2}+n}\right)\right),\\
&\chi_{1/2}(q)=\frac{1}{2}q^{-\frac{1}{48}}\left(\prod_{n=0}^\infty\left(1+q^{\frac{1}{2}+n}\right)-\prod_{n=0}^\infty\left(1-q^{\frac{1}{2}+n}\right)\right),\\
&\chi_{1/16}(q)= q^{\frac{1}{24}}\prod_{n=1}^\infty\left(1+q^n\right)
\end{align}


\subsubsection{Modular S Matrix and Transformation of Characters}
The modular S Matrix of Ising CFT is given by:
\begin{align}
  S^i_j=\begin{pmatrix}
    \frac{1}{2} & \frac{1}{2} & \frac{1}{\sqrt{2}} \\
    \frac{1}{2} & \frac{1}{2} & -\frac{1}{\sqrt{2}} \\
    \frac{1}{\sqrt{2}} & -\frac{1}{\sqrt{2}} & 0
  \end{pmatrix}
\end{align}
with the column and row order $ \mathbb{1},\epsilon,\sigma $. The characters form a representation of the modular group with the transformation:
\begin{align}
  \chi_i\left(q\right)=\sum_j S_{i}^j\chi_j\left(\tilde{q}\right).
\end{align}
if we take $ q = e^{-\frac{\pi \beta}{L}} $ and $ \tilde{q} = e^{-\frac{4\pi L}{\beta}} $, we have:
\begin{align}
  &\chi_0(\tilde{q})=\frac{1}{2}\left(\chi_0(q)+\chi_{1/2}(q)\right)+\frac{1}{\sqrt{2}}\chi_{1/16}(q),\\
  &\chi_{1/2}(\tilde{q})=\frac{1}{2}\left(\chi_0(q)+\chi_{1/2}(q)\right)-\frac{1}{\sqrt{2}}\chi_{1/16}(q),\\
  &\chi_{1/16}(\tilde{q})=\frac{1}{\sqrt{2}}\left(\chi_0(q)-\chi_{1/2}(q)\right).
\end{align}
as well as the inverse transformation:
\begin{align}
  &\chi_0(q)=\frac{1}{2}\left(\chi_0(\tilde{q})+\chi_{1/2}(\tilde{q})\right)+\frac{1}{\sqrt{2}}\chi_{1/16}(\tilde{q}),\\
  &\chi_{1/2}(q)=\frac{1}{2}\left(\chi_0(\tilde{q})+\chi_{1/2}(\tilde{q})\right)-\frac{1}{\sqrt{2}}\chi_{1/16}(\tilde{q}),\\
  &\chi_{1/16}(q)=\frac{1}{\sqrt{2}}\left(\chi_0(\tilde{q})-\chi_{1/2}(\tilde{q})\right).
\end{align}

\subsection{Ising BCFT}
For a BCFT induced from Ising CFT, we can construct its boundary states by solving the Ishibashi condition and the Cardy's condition.

\subsubsection{Ishibashi States of Ising BCFT}

The Ishibashi states of Ising CFT can be constructed from its three primary states:
\defi{
  \textbf{Ising Ishibashi States}

  Ising CFT has three Ishibashi states:
  \begin{align}
    &|0\rangle\rangle=\sum_{N}|0;N\rangle\otimes\overline{|0;N\rangle},\\
&|1/2\rangle\rangle=\sum_{N}|1/2;N\rangle\otimes\overline{|1/2;N\rangle},\\
&|1/16\rangle\rangle=\sum_{N}|1/16;N\rangle\otimes\overline{|1/16;N\rangle}.
  \end{align}
}

\subsubsection{Cardy States of Ising BCFT}

 Using the modular S matrix of Ising CFT, we can write down the Cardy states by definition:
\defi{
  \textbf{Ising Cardy States}

  Ising BCFT exist three Cardy states:
  \begin{align}
    &\left|+\right\rangle=\frac{1}{\sqrt{2}}|0\rangle\rangle+\frac{1}{\sqrt{2}}|1/2\rangle\rangle+\frac{1}{\sqrt[4]{2}}|1/16\rangle\rangle \\
&\left|-\right\rangle=\frac{1}{\sqrt{2}}|0\rangle\rangle+\frac{1}{\sqrt{2}}|1/2\rangle\rangle-\frac{1}{\sqrt[4]{2}}|1/16\rangle\rangle \\
& \left|f\right\rangle =|0\rangle\rangle- |1/2\rangle\rangle.
  \end{align}
}
In Cardy's original paper \cite{cardyBoundaryConditionsFusion1989}, these three boundary states are argued to correspond to the three physical boundary conditions of the Ising model\footnote{This is an argument made by considering the transformation of the boundary states under $ \mathbb{Z}_2 $ symmetry of Ising CFT. It is not a rigorous theorem.}:
\begin{itemize}
  \item $|+\rangle$: fixed boundary condition with spins fixed to $+1$ at the boundary.
  \item $|-\rangle$: fixed boundary condition with spins fixed to $-1$ at the boundary.
  \item $|f\rangle$: free boundary condition with spins free to fluctuate at the boundary.
\end{itemize}

\subsubsection{Partition Functions of Ising BCFT}

With the Cardy states, we can compute the partition functions of Ising BCFT on a compatified strip (cylinder) with different boundary conditions:
\begin{table}[H]
\centering
\renewcommand{\arraystretch}{1.6}
\begin{tabular}{cccc}
\hline\hline
Partition Function & Close Channel & Expression in $\tilde q$ & Expression in $q$ \\
\hline\hline

$Z_{++} = Z_{--}$ 
& $\langle \pm | \tilde{q}^{\frac12(L_0+\bar L_0 - c/12)} | \pm \rangle$
& $\frac12 \chi_0(\tilde q) + \frac12 \chi_{1/2}(\tilde q) + \frac{1}{\sqrt{2}} \chi_{1/16}(\tilde q)$
& $\chi_0(q)$
\\

$Z_{+-} = Z_{-+}$ 
& $\langle \pm | \tilde{q}^{\frac12(L_0+\bar L_0 - c/12)} | \mp \rangle$
& $\frac12 \chi_0(\tilde q) + \frac12 \chi_{1/2}(\tilde q) - \frac{1}{\sqrt{2}} \chi_{1/16}(\tilde q)$
& $\chi_{1/2}(q)$
\\

$Z_{+f} = Z_{-f}$
& $\langle \pm | \tilde{q}^{\frac12(L_0+\bar L_0 - c/12)} | f \rangle$
& $\frac{1}{\sqrt{2}} \chi_0(\tilde q) - \frac{1}{\sqrt{2}} \chi_{1/2}(\tilde q)$
& $\chi_{1/16}(q)$
\\

$Z_{ff}$
& $\langle f | \tilde{q}^{\frac12(L_0+\bar L_0 - c/12)} | f \rangle$
& $\chi_0(\tilde q) + \chi_{1/2}(\tilde q)$
& $\chi_0(q) + \chi_{1/2}(q)$
\\

\hline\hline
\end{tabular}
\caption{Partition function and corresponding characters in both $\tilde q$ and $q$ expansions.}
\label{tab:ising_partition_functions}
\end{table}

We can see that the expression in $ q $ always have non-negative integer coefficients, which satisfy the Cardy condition.


\subsection{Majorana CFT}
A Majorana CFT (also known as free massless Majorana fermion CFT) is a two-dimensional conformal field theory describing free massless Majorana fermions. I'll briefly review the canonical quantization of Majorana CFT here. 

\subsubsection{Definition of Majorana CFT}

A Majorana CFT is defined by the following action:
\defi{
  \textbf{Majorana Fermion Action}

  The action of a free Majorana fermion $\psi(z)$ and $\bar{\psi}(\bar{z})$ on a complex plane is given by:
  \begin{align}
    S = \frac{1}{2\pi} \int d^2z \, \left( \psi \bar{\partial} \psi + \bar{\psi} \partial \bar{\psi} \right),
  \end{align}
  where $\partial = \frac{\partial}{\partial z}$ and $\bar{\partial} = \frac{\partial}{\partial \bar{z}}$ and $ d^2 z = d^2 x = dx dy $. 
}
The equations of motion derived from the action are:
\begin{align}
  \bar{\partial} \psi(z) = 0, \quad \partial \bar{\psi}(\bar{z}) = 0,
\end{align}
which means that $\psi(z)$ is a holomorphic field and $\bar{\psi}(\bar{z})$ is an anti-holomorphic field. 

\subsubsection{Canonical Quantization}

To quantize the Majorana CFT, we impose radial quantization of the theory on a complex plane. The EoM gives us that the fields can be expanded as:
\begin{align}
  \psi(z)=\sum_kb_kz^{-k-1/2} \quad \bar{\psi}(\bar{z})=\sum_k\bar{b}_k\bar{z}^{-k-1/2}
\end{align}
it is a common convention to take the $ z^{1/2} $ in the mode expansion. The OPE of the fields gives us the anti-commutation relations of the modes:
\begin{align}
  \{ b_k, b_l \} = \delta_{k+l,0}, \quad \{ \bar{b}_k, \bar{b}_l \} = \delta_{k+l,0}, \quad \{ b_k, \bar{b}_l \} = 0.
\end{align}
Moreover the Majorana fermion condition $ \psi(z) = \psi^\dagger(z) $ gives us that the modes satisfy:
\begin{align}
  b_k^\dagger = b_{-k}, \quad \bar{b}_k^\dagger = \bar{b}_{-k}.
\end{align}
The Virasoro generators can be written in terms of the modes as:
\begin{align}
  L_n=\frac{1}{2}\sum_k(k+\frac{1}{2}):b_{n-k}b_k:
\end{align}

\subsubsection{Neveu-Schwarz and Ramond Sectors}

We can impose two different boundary conditions on the fermion fields:
\begin{itemize}
  \item \textbf{Neveu-Schwarz (NS) Boundary Condition}:
        \begin{align}
          \psi(e^{2\pi i} z) = \psi(z), \quad \bar{\psi}(e^{-2\pi i} \bar{z}) = \bar{\psi}(\bar{z}).
        \end{align}
        This boundary condition leads to half-integer moding of the fermion fields:
        \begin{align}
          k \in \mathbb{Z} + \frac{1}{2}.
        \end{align}
      \item \textbf{Ramond (R) Boundary Condition}: 
        \begin{align}
          \psi(e^{2\pi i} z) = - \psi(z), \quad \bar{\psi}(e^{-2\pi i} \bar{z}) =- \bar{\psi}(\bar{z}).
        \end{align}
        This boundary condition leads to integer moding of the fermion fields:
        \begin{align}
          k \in \mathbb{Z}.
        \end{align}
\end{itemize}
Different boundary conditions lead to different mode expansions of the fermion fields as well as different Hilbert spaces of the theory. The Hilbert space of the theory is thus classified into two sectors: the NS sector and the R sector.
The Hamiltonian of the Majorana CFT in both sectors can be written as:
\thm{
  \textbf{Majorana CFT Hamiltonian}

  The Hamiltonian of the Majorana CFT on a complex plane is given by:
  \begin{align}
    & H \;=\; L_0 + \bar L_0 \\
    %
    & L_0 \;=\; \sum_{k>0} k\, b_{-k} b_k ,
    \qquad 
      \bar L_0 \;=\; \sum_{k>0} k\, \bar b_{-k} \bar b_k ,
    \qquad (\mathrm{NS}:~ k\in \mathbb{Z}+\tfrac12) 
    \\
    & L_0 \;=\; \sum_{k>0} k\, b_{-k} b_k + \frac{1}{16},
    \qquad
      \bar L_0 \;=\; \sum_{k>0} k\, \bar b_{-k} \bar b_k + \frac{1}{16},
    \qquad (\mathrm{R}:~ k\in\mathbb{Z})
\end{align}
}

\subsubsection{Zero Mode and Ground State Degeneracy}

To construct the Hilbert space of the Majorana CFT, we need to define the ground states in both NS and R sectors. The NS sector works straightforwardly with a zero vacuum energy after normalization and a unique ground state $ \ket{0}_{NS} $ defined by:
\begin{align}
  b_k \ket{0}_{NS} = 0, \quad \bar{b}_k \ket{0}_{NS} = 0, \quad k > 0.
\end{align}
However, the R sector is more subtle due to the presence of zero modes $ b_0 $ and $ \bar{b}_0 $. These zero modes satisfy the Clifford algebra:
\begin{align}
  \{ b_0, b_0 \} = 1, \quad \{ \bar{b}_0, \bar{b}_0 \} = 1, \quad \{ b_0, \bar{b}_0 \} = 0.
\end{align}
We can see that the zero modes $ b_0 $ and $ \bar{b}_0 $ generate a two-dimensional representation of the Clifford algebra. Thus the ground state of the R sector is doubly degenerate according to the following theorm:
\thm{
  Two dimensional Cliffoed Algebra Representation have non-trivial representation for at least two dimensions.
}
We can define two ground states $ \ket{+}_{R} $ and $ \ket{-}_{R} $ and the act of pauli matrix on it is:
\begin{align}
  \sigma_x \ket{\pm}_{\mathrm R} &= \ket{\mp}_{\mathrm R}, \quad \sigma_y \ket{\pm}_{\mathrm R} = \mp i\,\ket{\mp}_{\mathrm R}, \quad \sigma_z \ket{\pm}_{\mathrm R} = \pm \ket{\pm}_{\mathrm R}.
\end{align}
And then we can represent the zero modes and the fermion parity operator as:
\begin{align}
   & b_0 = \frac{1}{2}\left( \sigma_x + \sigma_y \right)
        \,(-1)^{\sum_{n>0} b_{-n} b_n + \bar b_{-n} \bar b_n}, \\
   & \bar b_0 = \frac{1}{2}\left( \sigma_x - \sigma_y \right)
        \,(-1)^{\sum_{n>0} b_{-n} b_n + \bar b_{-n} \bar b_n}, \\
\end{align}
For the ground states, we have: 
\begin{align}\label{eq:Majorana zero mode action on R ground states}
  b_0|\pm\rangle_\mathrm{R}=\frac{1}{\sqrt{2}}e^{\pm i\pi/4}|\mp\rangle_\mathrm{R},\quad\bar{b}_0|\pm\rangle_\mathrm{R}=\frac{1}{\sqrt{2}}e^{\mp i\pi/4}|\mp\rangle_\mathrm{R}.
\end{align}

\subsubsection{Fermion Parity}\label{sec:Fermion Parity in Majorana CFT}

We hope to define (in many cases this operator might not be well defined) a fermion parity operator $ (-1)^F $ in a Fermionic CFT. 
\defi{
  \textbf{Fermion Parity Operator}

  This operator counts the parity of fermion number in the theory, satisfying:
\begin{align}
  \{(-1)^F, b_k\}=0, \quad \{(-1)^{F}, \bar{b}_k\}=0, \quad  ((-1)^F )^2 = 1.
\end{align}
}
In the NS sector, the fermion parity operator is explicitly given by:
\begin{align}
  (-1)^F = (-1)^{\sum_{k>0} b_{-k} b_k + \sum_{k>0} \bar{b}_{-k} \bar{b}_k}.
\end{align}
In the R sector, there are two different ways of defining the fermion number operator:
\begin{align}
  (-1)^{F} =
\begin{cases}
  -2 i\, b_0 \bar b_0 \ (-1)^{F_{nz}},  \\
  +2 i\, b_0 \bar b_0 \ (-1)^{F_{nz}}. 
\end{cases}
\end{align}
where $ (-1)^{F_{nz}} = (-1)^{\sum_{k>0} b_{-k} b_k + \sum_{k>0} \bar{b}_{-k} \bar{b}_k} $. In this note, we will use the first definition of fermion parity operator in the R sector.

\subsubsection{Hilbert Space Structure}

With the ground state structure and mode expansion, we can construct the Hilbert space of the Majorana CFT. The Hilbert space is classified into two sectors: the NS sector and the R sector. 
\begin{itemize}
  \item \textbf{NS Sector Hilbert Space}:
        The Hilbert space in the NS sector is constructed by acting the creation operators $ b_{-k}, \bar{b}_{-k} $ ($ k > 0 $ and $ k \in \mathbb{Z} + 1/2 $) on the unique ground state $ \ket{0}_{NS} $:
        \begin{align}
          \mathcal{H}_{NS} = \mathrm{span}\left\{ b_{-k_1} b_{-k_2} \cdots \bar{b}_{-l_1} \bar{b}_{-l_2} \cdots \ket{0}_{NS} \,|\, k_i, l_j > 0, k_i, l_j \in \mathbb{Z} + \frac{1}{2} \right\}.
        \end{align}
      \item \textbf{R Sector Hilbert Space}: 
        The Hilbert space in the R sector is constructed by acting the creation operators $ b_{-k}, \bar{b}_{-k} $ ($ k > 0 $ and $ k \in \mathbb{Z} $) on the doubly degenerate ground states $ \ket{\pm}_{R} $:
        \begin{align}
          \mathcal{H}_{R} = \mathrm{span}\left\{ b_{-k_1} b_{-k_2} \cdots \bar{b}_{-l_1} \bar{b}_{-l_2} \cdots \ket{\pm}_{R} \,|\, k_i, l_j > 0, k_i, l_j \in \mathbb{Z}  \right\}.
        \end{align}
\end{itemize}



